% Capítulo 1. Introducción
% Estructura según índice de la Dra. Lárraga.

\chapter{Introducción}
\label{ch:introduccion}

\begin{resumen}
Este capítulo presenta el contexto del crecimiento urbano acelerado en México y la selección de Querétaro como caso de estudio. Se identifican las limitaciones de los métodos tradicionales de modelado urbano y las brechas metodológica, temporal, contextual y operacional que motivan la investigación. Se justifica la integración de Weight of Evidence, Autómatas Celulares y Algoritmos Genéticos como solución innovadora. Se establecen los objetivos generales y específicos, las preguntas de investigación, las hipótesis, los alcances, las limitaciones y las contribuciones esperadas del trabajo.
\end{resumen}

\section{Contexto y problema}

El crecimiento urbano acelerado en México representa uno de los desafíos más significativos del siglo XXI. Entre 1950 y 2020, la población urbana mexicana creció de 43\% a 80\%, concentrándose inicialmente en megalópolis y posteriormente dispersándose hacia ciudades intermedias como Querétaro, León y Aguascalientes.

En este contexto, las herramientas de modelado y simulación se han convertido en un recurso valioso para comprender y predecir patrones de expansión urbana. Entre estas herramientas, los autómatas celulares (AC) han demostrado ser efectivos para representar la dinámica espacial de las ciudades debido a su capacidad para capturar interacciones locales y emergencias globales.

\subsection{Querétaro como caso de estudio}

Querétaro representa un caso paradigmático de crecimiento urbano acelerado en México central. La Tabla~\ref{tab:queretaro_facts} resume las características que justifican su selección.

\begin{table}[htbp]
\centering
\caption{Características de Querétaro como Caso de Estudio}
\label{tab:queretaro_facts}
\begin{tabular}{lp{9cm}}
\toprule
\textbf{Dimensión} & \textbf{Descripción} \\
\midrule
Crecimiento demográfico & De 456{,}000 habitantes (1990) a 1{,}049{,}777 (2020), representando un incremento del 130\% en tres décadas \\
Expansión territorial & Incremento de 340\% en área urbana durante el período 1984-2020 \\
Transformación económica & Transición de economía predominantemente agrícola a industrial-servicios \\
Posición estratégica & Ubicación privilegiada en el corredor del Bajío mexicano con conectividad multimodal \\
Atracción de inversión & Consolidación como hub industrial y tecnológico emergente a nivel nacional \\
\bottomrule
\end{tabular}
\end{table}

\subsection{Limitaciones de los enfoques tradicionales}

Los métodos convencionales de modelado urbano presentan limitaciones estructurales. Predominan los enfoques reactivos que responden a cambios ya ocurridos en lugar de anticiparlos. Los modelos empleados resultan excesivamente simplificados para capturar la complejidad de las dinámicas urbanas reales.

\begin{table}[htbp]
\centering
\caption{Limitaciones de Métodos Tradicionales de Modelado Urbano}
\label{tab:method_limitations}
\begin{tabular}{lp{10cm}}
\toprule
\textbf{Método} & \textbf{Limitaciones Principales} \\
\midrule
Modelos estadísticos clásicos & Asumen relaciones lineales, ignoran efectos espaciales y retroalimentación \\
Autómatas Celulares tradicionales & Emplean reglas de transición heurísticas con calibración manual subjetiva \\
Modelos de Machine Learning & Funcionan como ``caja negra'' con interpretabilidad limitada, requieren datos masivos \\
\bottomrule
\end{tabular}
\end{table}

\subsection{Brechas identificadas}

El análisis de literatura revela cuatro brechas críticas que motivan esta investigación. La \textit{brecha metodológica} evidencia la falta de integración sistemática entre evidencia empírica y simulación espacial. La \textit{brecha temporal} señala la validación limitada en horizontes de 20-30 años, período relevante para planificación urbana. La \textit{brecha contextual} indica la escasez de estudios específicos para ciudades intermedias mexicanas. Finalmente, la \textit{brecha operacional} refleja la carencia de herramientas accesibles para planificadores urbanos.

\section{Justificación}

La integración de técnicas complementarias —Weight of Evidence, Autómatas Celulares y Algoritmos Genéticos— ofrece una solución innovadora a estas limitaciones. WoE aporta cuantificación objetiva de relaciones espaciales. Los AC permiten simulación espacialmente explícita, representando patrones emergentes. Los AG aportan optimización automática de parámetros, logrando calibración objetiva y reproducible.

\section{Objetivos}

\subsection{Objetivo general}

Desarrollar un modelo evolutivo basado en autómatas celulares y algoritmos genéticos que permita simular y predecir el crecimiento urbano de la Zona Metropolitana de Querétaro (1984-2020), validado mediante el caso de estudio de Querétaro, que supere en precisión y robustez temporal a métodos de referencia.

\subsection{Objetivos específicos}

\begin{enumerate}
	\item Implementar un autómata celular que modele el cambio de uso de suelo en celdas urbanas y no urbanas, integrando evidencias empíricas calculadas mediante Weight of Evidence.
	\item Integrar un algoritmo genético para la optimización de parámetros del modelo, incluyendo pesos de influencia, umbrales y factores espaciales.
	\item Procesar y clasificar datos satelitales históricos (1984-2020) de Google Earth para construir mapas binarios de entrada al modelo.
	\item Validar el desempeño del modelo mediante métricas espaciales y temporales (IoU, Kappa, FoM), comparando las simulaciones con datos reales.
	\item Evaluar el impacto de diferentes configuraciones y escenarios sobre la predicción del crecimiento urbano.
\end{enumerate}

\section{Preguntas de investigación e hipótesis}

\subsection{Preguntas de investigación}

\begin{enumerate}
	\item ¿Es posible mejorar la precisión de un modelo de autómata celular para simular crecimiento urbano mediante optimización evolutiva con algoritmos genéticos?
	\item ¿Qué configuración de parámetros produce la mejor concordancia entre simulaciones y datos históricos?
	\item ¿Qué variables espaciales tienen mayor influencia en la expansión urbana de Querétaro?
	\item ¿En qué medida la integración de WoE, AC y AG mejora la robustez temporal de las predicciones en horizontes de 10-20 años?
\end{enumerate}

\subsection{Hipótesis}

\begin{itemize}
	\item \textbf{H1}: La metodología integrada WoE-AC-AG supera significativamente el desempeño de métodos de referencia (AC tradicional, regresión logística, extrapolación lineal) en precisión predictiva y robustez temporal, alcanzando métricas superiores (IoU $>$ 0.65, Kappa $>$ 0.60).
	\item \textbf{H2}: La calibración con AG mejora el fitness del modelo en al menos 15\% sobre calibración manual o heurística.
	\item \textbf{H3}: Los pesos WoE empíricos proporcionan mayor capacidad predictiva que pesos heurísticos basados en distancia.
	\item \textbf{H4}: El modelo mantiene capacidad predictiva estable (degradación $<$ 10\%) en validación temporal de 10 años.
\end{itemize}

\section{Alcances y limitaciones}

\subsection{Alcances}

\begin{itemize}
	\item El modelo propuesto se centra en la Zona Metropolitana de Querétaro, con un período de análisis de 37 años (1984-2020).
	\item Se utilizarán imágenes satelitales multitemporales de Google Earth para generar series históricas de uso de suelo.
	\item El sistema podrá simular escenarios futuros de expansión urbana a partir de condiciones iniciales definidas.
	\item La metodología está diseñada para ser transferible a otras ciudades intermedias mexicanas.
\end{itemize}

\begin{table}[htbp]
\centering
\caption{Delimitaciones y Alcances de la Investigación}
\label{tab:scope}
\begin{tabular}{llp{7cm}}
\toprule
\textbf{Dimensión} & \textbf{Aspecto} & \textbf{Especificación} \\
\midrule
\multirow{3}{*}{Espacial} & Área de estudio & Zona Metropolitana de Querétaro \\
 & Extensión & $\sim$1{,}200 km² incluyendo municipios conurbados \\
 & Resolución & $\sim$30 m/píxel (imágenes Landsat) \\
\midrule
\multirow{3}{*}{Temporal} & Período de análisis & 1984-2020 (37 años) \\
 & Calibración & 1984-2014 (30 años) \\
 & Validación & 2015-2020 (5 años) \\
\midrule
\multirow{3}{*}{Temática} & Enfoque & Crecimiento urbano horizontal (sprawl) \\
 & Representación & Clasificación binaria urbano/no-urbano \\
 & Variables & Factores de vecindad, densidad urbana, WoE espacial \\
\bottomrule
\end{tabular}
\end{table}

\subsection{Limitaciones}

\begin{itemize}
	\item La precisión del modelo depende de la calidad y resolución de los datos disponibles (imágenes RGB sin bandas infrarrojas).
	\item La representación binaria simplifica la complejidad urbana al no capturar densidades variables ni usos mixtos del suelo.
	\item Factores no espaciales, como políticas públicas futuras y crisis económicas, no son modelados explícitamente.
	\item La transferibilidad geográfica directa es limitada; se requiere recalibración para diferentes contextos urbanos.
\end{itemize}

\section{Contribuciones esperadas}

\begin{table}[htbp]
\centering
\caption{Contribuciones Esperadas de la Investigación}
\label{tab:contributions}
\begin{tabular}{lp{10cm}}
\toprule
\textbf{Dimensión} & \textbf{Contribuciones} \\
\midrule
Teórica & Framework metodológico para integración sistemática de WoE, AC y AG; protocolo estándar de validación temporal con series de 37 años \\
Metodológica & Algoritmo de calibración automática para parámetros AC; conjunto integral de métricas espaciales (IoU, Kappa, FoM, morfológicas) \\
Empírica & Dataset histórico 1984-2020 de Querétaro; caracterización de patrones de crecimiento urbano mexicano \\
Práctica & Pipeline reproducible open-source; directrices para planificación urbana sostenible; marco replicable para ciudades intermedias \\
\bottomrule
\end{tabular}
\end{table}

\section{Organización del documento}

El presente trabajo se organiza en seis capítulos:
\begin{itemize}
	\item \textbf{Capítulo 1 (Introducción):} Presenta el contexto, problema, objetivos, hipótesis, alcances y contribuciones.
	\item \textbf{Capítulo 2 (Marco teórico):} Revisa la literatura relevante sobre AC, WoE, AG y métricas de evaluación, justificando la importancia del estudio.
	\item \textbf{Capítulo 3 (Área de estudio):} Describe la obtención, preprocesamiento y análisis de los datos satelitales de Querétaro.
	\item \textbf{Capítulo 4 (Modelo de crecimiento urbano):} Detalla el diseño del autómata celular WoE-AC, la optimización con AG y el framework de validación.
	\item \textbf{Capítulo 5 (Resultados y análisis):} Presenta y analiza las simulaciones con métricas cuantitativas y comparaciones espaciales.
	\item \textbf{Capítulo 6 (Conclusiones):} Sintetiza hallazgos, discute implicaciones para planificación urbana y propone trabajo futuro.
\end{itemize}
